\section{Introduction}

Daian et al.~\cite{daian2020flash} introduced the term ``miner extractable value'' for the profit a miner can make by reordering, inserting or censoring transactions. After Ethereum moved to proof-of-stake the term was generalised to ``maximal extractable value'', since validators, block builders and independent searchers all take part in extraction~\cite{yang2022sok}.

MEV takes several forms. \emph{Arbitrage} closes price gaps between venues. \emph{Liquidation} captures the bonus paid for closing under-collateralised loans. \emph{Sandwich attacks} surround a user's swap with two trades to profit from its price impact~\cite{zhou2021high}. Qin et al.~\cite{qin2022quantifying} measured these strategies on Ethereum and found extracted value in the hundreds of millions of dollars. The root cause is the transparency of the mempool: pending transactions are visible to everyone before they execute~\cite{lyu2022empirical}.

\section{Constant Function Market Makers}

CFMMs are the dominant design for on-chain exchanges. A CFMM keeps reserves of two or more assets on an invariant curve $F(x,y)=k$; the constant product market maker sets $F(x,y)=x\cdot y$~\cite{angeris2020improved,adams2020uniswap}.

Two features of CFMMs matter for MEV. The first is \emph{order sensitivity}: running the same trades in a different order gives every party a different outcome, and this is what a sandwich attacker exploits. Yang et al.~\cite{yang2022sok} give a systematic account of the resulting opportunities and of the countermeasures proposed against them. The second is \emph{stale pricing}: the pool price moves only when someone trades, so arbitrageurs profit by trading against it whenever the external price changes. Milionis et al.~\cite{milionis2022lvr} formalise the resulting cost to liquidity providers as loss-versus-rebalancing (LVR). A redistribution mechanism should therefore compensate LPs as well as users.

\section{Private Channels}

The most widely deployed defence routes transactions directly to block builders, bypassing the public mempool~\cite{flashbots2022}. This hides a trade from outside searchers. It does not hide it from the relay or the builder, who see the full plaintext. Lyu et al.~\cite{lyu2022empirical} study private transactions on Ethereum and the security implications of this shift in trust. Proposer-builder separation (PBS) has a similar character: it limits what validators can extract, but concentrates block building among a small number of builders and returns nothing to users~\cite{heimbach2023pbs}.

\section{Batch Auctions}

Budish et al.~\cite{budish2015hft} proposed frequent batch auctions for traditional markets: orders are collected over a short interval and cleared together at a uniform price. Because every trade in a batch receives the same price, there is nothing to gain from being placed immediately before or after another trade. On-chain batch auction protocols apply the same idea. Their cost is latency, since a user must wait for the batch to close, and they do not redistribute MEV that arises between batches.

\section{Commit-Reveal Schemes and Encrypted Mempools}

A cryptographic commitment lets a party fix a value without disclosing it. With a hash commitment the party publishes $c=H(m\,\|\,r)$ for a random salt $r$ and later reveals $(m,r)$. The scheme is \emph{binding}, because finding a different opening would require a hash collision, and \emph{hiding}, because $c$ reveals nothing about $m$ when $r$ is random. Commit-reveal has been proposed as a front-running defence~\cite{yang2022sok}, but it requires the user to return and reveal, and a user who dislikes the outcome can simply refuse.

Encrypted mempools remove the second step. Transactions are encrypted under a public key whose secret key is shared among a committee, and are decrypted only after their position is fixed~\cite{blindperm2023}. The underlying tools are secret sharing~\cite{shamir1979share} and threshold cryptosystems~\cite{desmedt1989threshold}, in which any $t+1$ of $n$ parties can decrypt but $t$ cannot. Pedersen~\cite{pedersen1991threshold} showed how the committee can generate such a key jointly so that no trusted dealer ever holds it. Encrypted mempools prevent front-running but, on their own, do not return MEV to anyone.

\section{Order Fairness}

A different line of work constrains the ordering itself. Kelkar et al.~\cite{kelkar2020order} define order-fairness for Byzantine consensus: if enough nodes receive one transaction before another, the first must be ordered first. Cachin et al.~\cite{cachin2022quick} give a more efficient protocol with a related guarantee. These protocols remove the block producer's discretion. They are, however, sensitive to network latency, which a well-connected adversary can exploit, and they leave transaction contents visible.

\section{MEV Redistribution}

RediSwap~\cite{zhang2024rediswap} accepts that some MEV is unavoidable and redistributes it. Arbitrageurs report their price beliefs, each MEV opportunity is sold in a second-price auction~\cite{vickrey1961counterspeculation}, and the payment is refunded to users (for sandwich MEV) and LPs (for rebalancing MEV). The design is incentive compatible, and its sandwich construction makes opportunities independent of one another.

Its weakness is the transport layer. RediSwap receives user transactions through a private channel operated by a sequencer. An operator who is malicious or compromised can read transactions before the auction, extract the MEV directly, or pass the information to a favoured arbitrageur. The redistribution guarantee therefore holds only while the sequencer is honest.

\section{Comparison of Existing Approaches}

Table~\ref{tab:comparison} compares the approaches on four properties. The entries for existing approaches summarise the discussion above; the entry for MEVRES is justified in Chapters~\ref{ch:new} and~\ref{ch:result}.

\begin{table}[h]
\centering
\caption{Comparison of MEV Mitigation Approaches}
\label{tab:comparison}
\begin{tabular}{|L{3.6cm}|C{2.2cm}|C{2.4cm}|C{2.4cm}|C{2.4cm}|}
\hline
\textbf{Approach} & \textbf{Prevents front-running} & \textbf{Decentralized trust} & \textbf{MEV redistribution} & \textbf{Incentive compatible} \\
\hline
Private channels~\cite{flashbots2022} & Partial & No & No & N/A \\
\hline
Batch auctions~\cite{budish2015hft} & Yes & Yes & No & Partial \\
\hline
Commit-reveal~\cite{yang2022sok} & Partial & Yes & No & Partial \\
\hline
Order fairness~\cite{kelkar2020order,cachin2022quick} & Yes & Partial & No & N/A \\
\hline
Encrypted mempool~\cite{blindperm2023} & Yes & Yes & No & N/A \\
\hline
RediSwap~\cite{zhang2024rediswap} & Yes & No & Yes & Yes \\
\hline
\textbf{MEVRES (this work)} & \textbf{Yes} & \textbf{Yes} & \textbf{Yes} & \textbf{Yes} \\
\hline
\end{tabular}
\end{table}

\section{Summary}

Existing work offers either privacy without redistribution (encrypted mempools, order fairness, batch auctions) or redistribution without decentralized privacy (RediSwap). No approach in Table~\ref{tab:comparison} provides both. MEVRES, presented in the next chapter, combines the encrypted-mempool idea with RediSwap's auction to close this gap.
